\scanpage{63}
Pour simplifier, on supposera que $C$ est de \unsure{t.f.} sur $\mathbb Q_p$. De plus, on a un hom
\diagram{\[\boxed{\begin{tikzcd}
\pi\arrow[r,"\alpha"]&G(\unsure{\mathbb Q_\ell|\mathbb Q_p})
\end{tikzcd}}\]}
Si le foncteur \deleted{\unsure{$\Phi$}} $T_p$ (ou, ce qui revient au même, \deleted{\unsure{pl.}} le foncteur $\Delta$) est pl. fidèle, dans cet cas seulement) $G$ est l'enveloppe algébrique de $\alpha$.

De plus, le composé
\diagram{\[
\begin{tikzcd}[column sep=small]
C\arrow[r,no head]&\operatorname{Mod}_f(\pi,\mathbb Q_p)\arrow[r,no head,"{\otimes_{\mathbb Q_p}C}"]&\operatorname{Mod}_f(\pi,C)\arrow[r,no head]&\operatorname{Mod}_f(C)
\end{tikzcd}
\]}
est un foncteur \uline{bigradué}, \deleted{\unsure{i.e. on a un}} (gradué par degré total, et par décomposition de Hodge)

donc on a
\diagram{\[\boxed{\begin{tikzcd}
Gm_C\arrow[r,shift left=.6ex,"i_1"]\arrow[r,shift right=.6ex,"i_2"']&G_C
\end{tikzcd}}\]}
avec
\[
\unsure{i'|i}(\lambda)=i_1(\lambda)i_2(\lambda)=i(\lambda):\quad Gm_C\longrightarrow G_C
\]
l'hom \uline{central} définissant la graduation de $C$.

De plus, on a un \deleted{hom} foncteur fibre \inserted{$T_{\mathrm{Hdg}}$} bigradué ``de Hodge'' sur $K$, donc défini par
\[
\left\{\begin{gathered}
\boxed{P,\ \text{torseur pour }G_L}\\
(\overline T_{\mathrm{Hdg}}=P\mathbin{\wedge}^{G_L}(T_p\otimes_{\mathbb Q_p}L)\ \deleted{\unsure{et}}
\end{gathered}\right.
\]
\deleted{et}
\diagram{\[
\left\{\begin{gathered}
\begin{tikzcd}
Gm_L\arrow[r,shift left=.6ex,"j_1"]\arrow[r,shift right=.6ex,"j_2"']&\deleted{\unsure{G'}}G_L^P=\operatorname{ad}(P)
\end{tikzcd}\\
j_1(\lambda)j_2(\lambda)=i(\lambda)
\end{gathered}\right.
\]}
\deleted{\unsure{$T_{\mathrm{Hdg}}$}}

avec un isom
\mspage{62}
